%------------------------
% Ethan's Résumé Template
% Author: necusjz
% License: CC-BY-4.0
%------------------------
\documentclass[letterpaper,11pt]{article}

\usepackage{latexsym}
\usepackage[empty]{fullpage}
\usepackage{titlesec}
\usepackage{marvosym}
\usepackage[usenames,dvipsnames]{color}
\usepackage{verbatim}
\usepackage{enumitem}
\usepackage[hidelinks]{hyperref}
\usepackage{fancyhdr}
\usepackage[english]{babel}
\usepackage{tabularx}
\usepackage{fontawesome5}
\usepackage{ragged2e}
\usepackage{etoolbox}
\usepackage{tikz}
\input{glyphtounicode}

% font options
\usepackage{newpxtext}
\linespread{1.05}  % palladio needs more leading (space between lines)
\usepackage[T1]{fontenc}

\pagestyle{fancy}
\fancyhf{}  % clear all header and footer fields
\fancyfoot{}
\renewcommand{\headrulewidth}{0pt}
\renewcommand{\footrulewidth}{0pt}

% adjust margins
\addtolength{\oddsidemargin}{-0.5in}
\addtolength{\evensidemargin}{-0.5in}
\addtolength{\textwidth}{1in}
\addtolength{\topmargin}{-.5in}
\addtolength{\textheight}{1.0in}

\urlstyle{same}

\raggedbottom
\raggedright
\setlength{\tabcolsep}{0in}
\setlength{\footskip}{5pt}

% sections formatting
\titleformat{\section}{
  \vspace{-2pt}\scshape\raggedright\large
}{}{0em}{}[\color{black}\titlerule\vspace{-5pt}]

% ensure that generate pdf is machine readable/ATS parsable
\pdfgentounicode=1

% custom commands
\newcommand{\cvitem}[1]{
  \item\small{
    {#1\vspace{-2pt}}
  }
}

\newcommand{\cvheading}[4]{
  \vspace{-2pt}\item
    \begin{tabular*}{\textwidth}[t]{l@{\extracolsep{\fill}}r}
      \textbf{#1} & #2 \\
      \small#3 & \small #4 \\
    \end{tabular*}\vspace{-7pt}
}

\newcommand{\cvheadingstart}{\begin{itemize}[leftmargin=0in, label={}]}
\newcommand{\cvheadingend}{\end{itemize}}
\newcommand{\cvitemstart}{\begin{itemize}[label=\textopenbullet]\justifying}
\newcommand{\cvitemend}{\end{itemize}\vspace{-5pt}}

\newcommand{\cvskill}[2]{
  \textcolor{black}{\textbf{#1}}\hfill
  \foreach \x in {1,...,5}{%
    \space{\ifnumgreater{\x}{#2}{\color{black!80!white!20}}{\color{black}}\faSquare}}\par%
  \vspace{-2pt}
}

\renewcommand\labelitemii{$\vcenter{\hbox{\footnotesize$\bullet$}}$}

\begin{document}

% contact information
\begin{center}
  \textbf{\LARGE\scshape Yunseong Kim} \\
  \vspace{1pt}\small
  \href{mailto:p4ranlee@gmail.com}{\textbf{Email}}
  $\ \diamond\ $ ysk@kzalloc.com
  \textbf{Phone}
  $\ \diamond\ $+46 70-260 30 51
  \href{https://github.com/yskzalloc}{\textbf{GitHub}}
  \href{https://www.linkedin.com/in/yunseong-kim-linux-kernel/}{\textbf{LinkedIn}}
\end{center}

\section{Professional Experience}

\cvheadingstart
  \cvheading
    {Ericsson Software Technology}{Stockholm, Sweden}
    {Senior Linux Engineer}{04/2026 - Present}
  \cvitemstart
    \cvitem{Linux Kernel Opensource Team}
  \cvitemend
  
  \cvheading
    {Ericsson Korea}{Seoul, South Korea}
    {System Software Developer}{07/2023 - 04/2026}
  \cvitemstart
    \cvitem{Mobility and Control Development Team}
  \cvitemend

  \cvheading
    {AhnLab}{Bundang-gu, South Korea}
    {System Software Developer}{01/2022 - 06/2023}
  \cvitemstart
    \cvitem{EDR(Endpoint Detection and Response) Development Team}
  \cvitemend

  \cvheading
    {Innowireless}{Bundang-gu, South Korea}
    {Software Developer}{8/2018 - 12/2021}
  \cvitemstart
    \cvitem{Server Software Development Team}
\cvitemend

\cvheadingend

\section{Education}
\cvheadingstart
  \cvheading
    {Korea Aerospace University}{Goyang, South Korea}
    {B.S. in Electronics and Avionics Engineering}{03/2010 - 02/2018}
  \cvheading
    {Linux Foundation}{Online}
    {SKF201: Threat Modeling}{9/2025 - 9/2025}
  \cvheading
    {Linux Foundation}{Online}
    {LFD441: Security and the Linux Kernel}{7/2024 - 7/2024}
\cvheadingend

\section{Key Performances}
\cvheadingstart
\item
\cvitemstart
  \cvitem{\textbf{Linux Kernel Security Development}}
  \begin{itemize}
    \item Developed security modules using \textbf{LSM}, \textbf{kprobe}, and \textbf{eBPF}
  \end{itemize}

\cvitem{\textbf{\href{https://lore.kernel.org/all/?q=Yunseong}{Upstream Linux Kernel Contributions}}}
  \begin{itemize}
    \item Reported 4 security vulnerabilities: merged patches for the mainline kernel: \href{https://lore.kernel.org/linux-cve-announce/2025081651-CVE-2025-38510-f67d@gregkh/}{\#1}, \href{https://lore.kernel.org/linux-cve-announce/2024122721-CVE-2024-53186-7c05@gregkh/}{\#2}, \href{https://lore.kernel.org/linux-cve-announce/2024080740-CVE-2024-42235-efae@gregkh/}{\#3}, \href{https://lore.kernel.org/linux-cve-announce/2024072929-CVE-2024-41021-f857@gregkh/}{\#4}
  \end{itemize}

 \cvitem{\href{http://perf.wiki.kernel.org/}{\textbf{Maintainer of the Linux perf tools wiki}};} {\textbf{uftrace contributions}: \href{https://github.com/namhyung/uftrace/pulls?q=author\%3Aparanlee}{\#1}, \href{https://github.com/namhyung/uftrace/pulls?q=author%3Akzall0c+}{\#2}}

  \cvitem{\textbf{Security R\&D in 5G Core Systems}}
  \begin{itemize}
    \item Developing \textbf{eBPF + rust} (aya-rs) based NGAP firewall
    \item Find vulnerabilities from \textbf{ASN.1 codec} to \textbf{5G NGAP protocol handling}
  \end{itemize}

\cvitemend
\cvheadingend

\section{Technical Skills}
\cvheadingstart
\item
\cvitemstart
    \cvitem{\textbf{Linux Kernel Security Module} Development using QEMU, LSM, crash tools, ftrace, perf tools and Wireshark}
    \cvitem{Strong C/C++ programming and debugging skills}
    \cvitem{\textbf{Fuzzing} with \textbf{syzkaller} and \textbf{AFL++}; source control with Git; Jenkins; scripting with Erlang/OTP, Python, and shell}
    \cvitem{\textbf{Linux Kernel CI/CD and Verification} Automation}
    \begin{itemize}
      \item Managed \textbf{multiple Jenkins pipelines} on virtualized infrastructures: VMware ESXi, vSphere
    \end{itemize}

\cvitemend
\cvheadingend

\section{Talks \& Publications}
\cvheadingstart
\item
\cvitemstart

\cvitem{\href{https://youtu.be/__DURK3p218?si=aPn0f2Kixq3V4P7T}{{\textbf{Speaker at DebConf25:}: \textit{Rethinking Cryptography in the Linux Kernel – Preparing for the Post-Quantum (PQC) Era}}} }
 \cvitem{
  \href{https://youtu.be/nGxzzORZm1o?si=1k9Mz54uqYNatnO1}{\textbf{Head First Reporting of Linux Kernel CVEs:} \textit{Practical Use of the Kernel Fuzzer}}}
\begin{itemize}
  \item Presented Linux kernel CVE discovery using syzkaller, including KSMBD and s390x vulnerabilities
\end{itemize}

  \cvitem{\href{https://youtu.be/RRJN8JS5xAE?si=_X-7-OyXzPrcyGes}{{\textbf{Speaker at DebConf24:}} \textit{perf In Action}} – page fault  comparison and fuzzing use cases and techniques for Linux perf}

  \cvitem{ \textbf{\href{https://www.linkedin.com/posts/yunseong-kim-linux-kernel_xilinx-fpga-linux-activity-7341153544160362496-_sBg?utm_source=share&utm_medium=member_desktop&rcm=ACoAADZ7vgEBrd70u0-SnyUUHL-swTgyhcvNQGE}{ZYBO Z7-20 AES Crypto Accelerator}} HW/SW co-development on the \textit{\textbf{PetaLinux v2025.1}}}

  \cvitem{\href{https://www.hackster.io/p4ranlee/zybo-z7-20-webcam-on-the-petalinux-c28709}{{\textbf{Hackster.io:}} \textit{ZYBO Z7-20 Webcam on the PetaLinux Project}} }
  
\cvitemend
\cvheadingend


\end{document}
